\section{Theoretical foundation and algorithm}

\subsection{Cubic-equation model used in the design}

The hardware is designed to solve the general cubic equation
\begin{equation}
ax^3 + bx^2 + cx + d = 0
\end{equation}
with all coefficients represented in IEEE-754 single-precision floating-point format. Instead of solving the equation directly in its original form, the design first standardizes the coefficients by dividing them by $-3a$:
\begin{equation}
\hat{b} = \frac{b}{-3a}, \qquad
\hat{c} = \frac{c}{-3a}, \qquad
\hat{d} = \frac{d}{-3a}
\end{equation}
In the source code and in the golden model, these normalized quantities are stored back into the working registers \texttt{r\_b}, \texttt{r\_c}, and \texttt{r\_d}.

After standardization, the solver evaluates the two characteristic quantities
\begin{equation}
\Delta_0 = \hat{b}^2 + \hat{c}
\end{equation}
\begin{equation}
\Delta_1 = 2\hat{b}^3 + 3\hat{b}\hat{c} + 3\hat{d}
\end{equation}
and then computes the discriminant-like term
\begin{equation}
\Delta = \Delta_1^2 - 4\Delta_0^3
\end{equation}
These three values determine which branch of the solution formula must be used.

\subsection{Root-generation branches}

The implemented algorithm follows the Cardano-style decomposition summarized in the mathematical notes:
\begin{itemize}
    \item If $\Delta < 0$, the square root of $\Delta$ is purely imaginary. The solver therefore computes $\sqrt{-\Delta}$ on the real datapath, treats the result as the imaginary part of $\sqrt{\Delta}$, and then performs a complex cube-root iteration to obtain $C$.
    \item If $\Delta \ge 0$ and the special repeated-root condition is detected, the design uses a simplified real cube-root computation based on $\Delta_1$.
    \item In the remaining case, the design computes the real quantity
    \begin{equation}
        C = \sqrt[3]{\frac{\Delta_1 + \sqrt{\Delta}}{2}}
    \end{equation}
    and reconstructs the three roots from $C$ and $\Delta_0/C$.
\end{itemize}

Let
\begin{equation}
\varepsilon = -\frac{1}{2} + j\frac{\sqrt{3}}{2}
\end{equation}
be the primitive cubic root of unity. The final roots are generated by one of the following forms:
\begin{equation}
x_k = \hat{b} + \varepsilon^k C + \varepsilon^{2k}\overline{C}, \qquad k \in \{0,1,2\}
\end{equation}
for the complex-conjugate branch, or
\begin{equation}
x_k = \hat{b} + \varepsilon^k C
\end{equation}
for the simplified repeated-root branch, or
\begin{equation}
x_k = \hat{b} + \varepsilon^k C + \varepsilon^{2k}\frac{\Delta_0}{C}
\end{equation}
for the general real branch.

\subsection{Pad\'e iterative approximation for square root}

The square-root block inside the solver does not use a dedicated CORDIC or LUT-based implementation. Instead, it relies on the Pad\'e iterative form already validated in the mathematical model:
\begin{equation}
P_{n+1} = P_n \cdot \frac{P_n^2 + 3S}{3P_n^2 + S}
\end{equation}
where $S$ is the input operand and $P_n$ is the approximation at iteration $n$.

This form is especially convenient for the current architecture because every iteration can be decomposed into a short sequence of floating-point multiplications, additions, and one division. The initial estimate is derived directly from the exponent field of the input FP32 number, which reduces the number of iterations required to reach a stable result.

In hardware, the square-root controller stores:
\begin{itemize}
    \item \texttt{r\_sqrt\_input}: the operand $S$,
    \item \texttt{r\_sqrt\_output}: the current estimate $P_n$,
    \item \texttt{iteration\_count}: number of Pad\'e updates already performed.
\end{itemize}

\subsection{Pad\'e iterative approximation for cube root}

The cube-root block follows a similar principle. The Pad\'e update used in both the software model and the Verilog implementation is
\begin{equation}
P_{n+1} = P_n \cdot \frac{P_n^3 + 2S}{2P_n^3 + S}
\end{equation}
This recurrence is used in two different modes:
\begin{itemize}
    \item \textbf{Real mode}: used when the cube-root argument is purely real. In this case only \texttt{r\_cbrt\_re\_input} and \texttt{r\_cbrt\_re\_output} are active.
    \item \textbf{Complex mode}: used when the argument is complex. The datapath then tracks both real and imaginary parts separately using \texttt{r\_cbrt\_re\_input}, \texttt{r\_cbrt\_im\_input}, \texttt{r\_cbrt\_re\_output}, and \texttt{r\_cbrt\_im\_output}.
\end{itemize}

Because the solver is implemented entirely with FP32 operators, the initial cube-root estimate is also built from the exponent field. A small helper block, \texttt{Divide\_int}, approximates the exponent scaling required for $\sqrt[3]{S}$ and avoids introducing a separate floating-point pre-processing unit.

\subsection{The golden model}

Before building the Verilog architecture, the algorithm was first implemented in the notebook \texttt{golden\_model.ipynb}. This notebook plays three important roles in the project:
\begin{itemize}
    \item It defines a software reference for \texttt{pade\_sqrt()}, \texttt{pade\_cbrt()}, and \texttt{solve\_cubic()} so that the team can validate the mathematical flow independently from hardware timing issues.
    \item It benchmarks the iterative square-root and cube-root formulas on large randomized datasets to evaluate convergence and numerical error.
    \item It generates random FP32 hexadecimal test vectors and compares the simulated Verilog outputs against the expected roots, thereby providing a complete reference-based verification loop.
\end{itemize}

Therefore, the golden model is not merely a prototype. It is the behavioral specification that guided the partitioning of the hardware datapath, the definition of the branch conditions, and the scheduling of arithmetic operations inside the finite state machine.
